\chapter{分场景案例分析}

\section{案例一：理想数据场景}
在理想数据下，参数恢复误差主要受积分截断与迭代停止阈值影响。实验显示当阈值从 $10^{-6}$ 收紧至 $10^{-8}$ 时，参数误差下降但时间开销增加明显。

\section{案例二：中噪声场景}
在中噪声水平（$\sigma\approx0.02$）下，高斯牛顿法能快速降残差，但后期易出现平台。引入阻尼后，平台区间缩短，整体误差更低。

\section{案例三：高噪声与偏差初值}
该场景最具挑战。多初值并行策略能显著提高可用解比例；若仅使用单初值，失败率显著升高。

\section{案例四：数据缺失窗口}
当观测存在间断时，建议采用分段目标函数与窗口权重，避免缺失段对全局估计造成不合理牵引。

\section{综合结论}
实际应用中应根据噪声水平与初值质量动态选择策略：低噪声优先速度，中高噪声优先稳定。
